\paragraph{Related work.} Our work was inspired by the works of 
\cite{YangCA24,barcelo2024logical,tale,LC25}, which exhibit a close connection
between unique-hard attention transformers and star-free regular languages.
In particular, these works also exploited the connection to LTL. However,
none of these results investigated neither the issue of succinctness nor
computational complexity of verification, which we establish in this
paper.

\cite{SAL25} investigated the issue of verifying transformers of various
precisions. In particular, it was shown that transformers of fixed precision
are at least $\NEXP$-hard (i.e. hard for the class of problems solvable by
nondeterministic algorithms that run in exponential time). The technique there
implies that transformers are (singly) exponentially more succinct than finite
automata. However, no conclusion can be derived as to their succinctness
relative to representations like LTL or RNN. Our result substantially improve
this by showing that transformers are doubly exponentially more succinct than 
automata, and exponentially more succinct than LTL and RNN. In addition, our
work assumes a much simpler model in comparison to the results in \citep{SAL25}. 
In particular, we use unique-hard attention, 
whereas in \citep{SAL25} a combination of softmax and hardmax is employed.
Finally, our results use positional masking (as employed in
\citep{YangCA24,tale,LC25}) --- as a simple class of Positional Embeddings 
(PEs) --- in contrast to \citep{SAL25}, which admits arbitrary PEs of fixed precision.

Succinctness has also been studied in the context of linguistics. For example,
according to Zipf's law of abbreviation \citep{zipf35}, frequently occuring
concepts tend to have a succinct description. In particular, Hindu-Arabic
numeral system --- which evolves into our modern numeral system --- allows
an exponentially more succinct description than the Roman numeral system.
According to Zipf's law, the former potentially enables mathematics and computer
science as we see today.
